% ---------------------------------------------------------------------------
% Preâmbulo: pacotes, aspeto e macros do relatório.
% ---------------------------------------------------------------------------

\usepackage[T1]{fontenc}
\usepackage[utf8]{inputenc}
\usepackage[portuguese]{babel}
\usepackage{lmodern}
\usepackage{microtype}
\usepackage[a4paper, margin=2.5cm, headheight=14pt]{geometry}
\usepackage{setspace}
\onehalfspacing

% Matemática
\usepackage{amsmath, amssymb, mathtools, bm}

% Tabelas
\usepackage{array, booktabs, tabularx, longtable, multirow, makecell}
\usepackage{colortbl}
\usepackage{threeparttable}
\usepackage{pdflscape}

% Figuras
\usepackage{graphicx}
\usepackage{float}
\usepackage[section]{placeins}  % figuras e tabelas não passam para a secção seguinte
\usepackage[font=small, labelfont=bf, labelsep=period]{caption}
\usepackage{subcaption}
\usepackage{tikz}
\usetikzlibrary{arrows.meta, positioning, shapes.geometric, shapes.multipart,
                fit, calc, backgrounds, decorations.pathreplacing}
\usepackage{pgfplots}
\pgfplotsset{compat=1.16}

% Diversos
\usepackage{xcolor}
\usepackage[inline]{enumitem}
\usepackage{pifont}
\usepackage{csquotes}
\usepackage{url}
\usepackage{fancyhdr}
\usepackage{titlesec}
\usepackage[hidelinks, unicode, pdfusetitle]{hyperref}

% ---------------------------------------------------------------------------
% Cores (inspiradas na identidade do Iscte; ajustáveis)
% ---------------------------------------------------------------------------
\definecolor{iscte}{RGB}{0, 51, 102}
\definecolor{acento}{RGB}{0, 112, 192}
\definecolor{cinza}{RGB}{110, 110, 110}
\definecolor{fundotab}{RGB}{235, 241, 248}
\definecolor{bom}{RGB}{0, 128, 0}
\definecolor{mau}{RGB}{190, 30, 45}
\definecolor{parcial}{RGB}{200, 130, 0}

% ---------------------------------------------------------------------------
% Notação UML e BPMN (estilos TikZ partilhados pelas figuras)
% ---------------------------------------------------------------------------
\tikzset{
  % UML: pacote (pasta com separador); o nome vai no corpo
  umlpkg/.style={draw=iscte, fill=fundotab, minimum height=7mm, minimum width=22mm, font=\small\ttfamily,
    align=center, append after command={\pgfextra{\draw[iscte, fill=fundotab]
      ([yshift=-0.1pt]\tikzlastnode.north west) rectangle ++(6mm,1.6mm);}}},
  % UML: pacote que contém outros (nome no separador)
  umlpkgbox/.style={draw=iscte, rounded corners=0pt, inner sep=7pt, append after command={\pgfextra{
      \node[draw=iscte, fill=white, anchor=south west, font=\small\bfseries\ttfamily, inner sep=2pt,
        minimum height=4mm] at ([yshift=-0.1pt]\tikzlastnode.north west) {#1};}}},
  % UML: componente (ícone no canto superior direito)
  umlcomp/.style={draw=iscte, fill=fundotab, minimum height=9mm, minimum width=22mm, font=\small, align=center,
    inner sep=3pt, append after command={\pgfextra{
      \draw[iscte, fill=white] ([xshift=-4.2mm, yshift=-1.2mm]\tikzlastnode.north east) rectangle ++(3mm,-2.6mm);
      \draw[iscte, fill=white] ([xshift=-5mm, yshift=-1.7mm]\tikzlastnode.north east) rectangle ++(1.6mm,-0.6mm);
      \draw[iscte, fill=white] ([xshift=-5mm, yshift=-2.9mm]\tikzlastnode.north east) rectangle ++(1.6mm,-0.6mm);}}},
  % UML: artefacto/base de dados (estereótipo no texto)
  umlart/.style={draw=iscte, fill=white, minimum height=8mm, minimum width=20mm, font=\small, align=center},
  % UML: nó de implantação (caixa 3D); o nome é dado por label
  umlnode/.style={draw=cinza, inner sep=7pt, append after command={\pgfextra{
      \draw[cinza] (\tikzlastnode.north west) -- ++(2mm,2mm) -| ([xshift=2mm, yshift=2mm]\tikzlastnode.south east)
        -- (\tikzlastnode.south east);
      \draw[cinza] (\tikzlastnode.north east) -- ++(2mm,2mm);}}},
  % UML: classe ou interface com compartimentos
  umlclass/.style={draw=iscte, fill=fundotab, rectangle split, rectangle split parts=2, font=\scriptsize\ttfamily,
    align=left, inner sep=3pt},
  % UML: setas
  umldep/.style={-{Straight Barb[length=2mm]}, dashed, thick, iscte},           % dependência («use»)
  umlreal/.style={-{Triangle[open, length=2.5mm, width=2.5mm]}, dashed, thick, iscte}, % realização
  umlassoc/.style={-{Straight Barb[length=2mm]}, thick, iscte},                % associação navegável / mensagem
  umlmsg/.style={-{Stealth[length=2mm]}, thick, iscte},                        % mensagem síncrona
  umlret/.style={-{Straight Barb[length=2mm]}, dashed, iscte},                  % resposta
  umlconn/.style={thick, iscte},                                               % conector/caminho de comunicação
  % UML (atividade) e BPMN
  bpmntask/.style={draw=iscte, fill=fundotab, rounded corners=3pt, minimum height=9mm, minimum width=19mm,
    text width=18mm, font=\scriptsize, align=center, inner sep=2pt},
  bpmnstart/.style={circle, draw=iscte, thick, fill=white, minimum size=5mm, inner sep=0pt},
  bpmnend/.style={circle, draw=iscte, line width=1.8pt, fill=white, minimum size=5mm, inner sep=0pt},
  bpmngw/.style={diamond, draw=iscte, thick, fill=white, minimum size=7mm, inner sep=0pt, font=\small\bfseries},
  bpmnflow/.style={-{Stealth[length=2mm]}, thick, iscte},
  actstart/.style={circle, fill=iscte, minimum size=3.5mm, inner sep=0pt},
  actend/.style={circle, draw=iscte, thick, fill=white, minimum size=4.5mm, inner sep=0pt,
    path picture={\fill[iscte] (path picture bounding box.center) circle (1.2mm);}},
}

% ---------------------------------------------------------------------------
% Títulos de secção e cabeçalhos
% ---------------------------------------------------------------------------
\titleformat{\section}{\Large\bfseries\color{iscte}}{\thesection}{0.8em}{}
\titleformat{\subsection}{\large\bfseries\color{iscte}}{\thesubsection}{0.8em}{}
\titleformat{\subsubsection}{\normalsize\bfseries\color{iscte}}{\thesubsubsection}{0.8em}{}

\pagestyle{fancy}
\fancyhf{}
\fancyhead[L]{\small\color{cinza}\TituloCurto}
\fancyhead[R]{\small\color{cinza}ADS --- Projeto}
\fancyfoot[C]{\small\thepage}
\renewcommand{\headrulewidth}{0.4pt}

% Nomes em português de Portugal
\addto\captionsportuguese{%
  \renewcommand{\contentsname}{Índice}%
  \renewcommand{\listfigurename}{Lista de Figuras}%
  \renewcommand{\listtablename}{Lista de Tabelas}%
  \renewcommand{\refname}{Referências}%
  \renewcommand{\abstractname}{Resumo}%
  \renewcommand{\tablename}{Tabela}%
  \renewcommand{\figurename}{Figura}%
  \renewcommand{\appendixname}{Anexo}%
}

% ---------------------------------------------------------------------------
% Macros
% ---------------------------------------------------------------------------
% Marcas usadas nas tabelas comparativas
\newcommand{\Sim}{\textcolor{bom}{\ding{51}}}          % abordado
\newcommand{\Parc}{\textcolor{parcial}{\ding{108}}}    % abordado parcialmente
\newcommand{\Nr}{\textcolor{cinza}{--}}                % não reportado
\newcommand{\Nao}{\textcolor{mau}{\ding{55}}}          % ausente por desenho


% Colunas de tabela
\newcolumntype{L}[1]{>{\raggedright\arraybackslash}p{#1}}
\newcolumntype{C}[1]{>{\centering\arraybackslash}p{#1}}
\newcolumntype{Y}{>{\raggedright\arraybackslash}X}

% Caixas de definição/observação
\newcommand{\caixa}[2]{%
  \begin{center}\begin{tikzpicture}
    \node[draw=acento, fill=fundotab, rounded corners=2pt, inner sep=8pt,
          text width=0.9\linewidth, align=justify] {\textbf{#1.}\ #2};
  \end{tikzpicture}\end{center}}

% ---------------------------------------------------------------------------
% Macros específicas do relatório de ADS
% ---------------------------------------------------------------------------
% Texto ainda por escrever: aparece a vermelho para não passar despercebido.
\newcommand{\porescrever}[1]{\textcolor{mau}{\itshape [Por escrever: #1]}}

% Inclui uma tabela/figura gerada a partir das medições (make dados).
% Se o ficheiro ainda não existir, mostra uma caixa com o nome esperado.
\newcommand{\gerado}[1]{%
  \IfFileExists{gerado/#1}{\input{gerado/#1}}%
    {\caixa{Dados em falta}{\texttt{gerado/\detokenize{#1}} --- correr \texttt{make dados}.}}}

% Identificadores de cenários, decisões e lacunas (ligam ao plano e aos ADRs).
\newcommand{\cen}[1]{\textsf{#1}}
\newcommand{\dec}[1]{\textsf{D#1}}
